\section{民事主体}
\begin{enumerate}
  \item 关键词解题法。识别 \yi 民事法律关系。财产、综合、人身；\er \zd{标的物}。不动产、权利、动产。
  \item 子女。婚生、非婚生、养子女、有抚养关系的继子女和胎儿。\al{家有儿女}。
  \item 自然人。
        \begin{enumerate}
          \item 民事权利\zd{能力}。做人的\zd{资格}，\zd{始于出生，终于死亡}。活着 \yjt 权利；生前、死后 \yjt 民事利益。
                \begin{enumerate}
                  \item 生前。胎儿利益的保护，1 个原则（不是人），3 个例外（\zd{纯获法律上的利益}）。
                        \begin{enumerate}
                          \item 遗产继承。遗腹子，胎儿预留份制度。\al{怀孕 8 个月，丈夫去世}。
                          \item 接受赠予。\al{能生下来送 10 根金条给孩子}。
                          \item 侵权行为。可提前行权。\al{撞了孕妇早产}。
                          \item \zd{娩出为死体，其民事权利能力自始不存在}。
                        \end{enumerate}
                  \item 死后。死者人格利益保护，\zd{7 项人格利益，8 类近亲属，2 个顺位，782}。
                        \begin{enumerate}
                          \item 姓名、肖像、\zd{名誉、荣誉}、隐私、遗体、遗骨 \yjt 保护的是近亲属情感利益。
                          \item \zd{配偶、父母、子女} | 兄弟姐妹、祖父母、外祖父母、孙子女、外孙子女。
                          \item 英雄烈士 \yjt 检察院公益诉讼。参看民诉、刑诉。
                        \end{enumerate}
                \end{enumerate}
          \item 民事行为能力。独立做事的资格。认识、识别、理解 \yjt 意思能力 = 年龄 + 精神状况。
                \begin{enumerate}
                  \item 无民事行为能力 < 8，一律无效。\al{去打酱油}。
                  \item 限制民事行为能力 < 18，不相适应的单方/双方民事法律行为 \yjt 无效/待定。
                  \item 完全民事行为能力 > 18 \hz\, \zd{16 + 劳动收入}。
                \end{enumerate}
          \item 监护。小的和老的。监护人（扩大解释，可以多个，可以是法人）。
                \begin{enumerate}
                  \item 确定。
                        \begin{enumerate}
                          \item 法定监护。第 1 顺位监护人\zd{法定性、首要性、当然性}。
                                \begin{itemize}
                                  \item 小的。\zd{父母}；祖父母、外祖父母；兄弟姐妹；其他个人 \hz 组织（愿意 + 同意）；单位监护人（居委会、村委会、民政部门）。
                                  \item 老的。配偶；父母、子女；其他近亲属；其他个人 \hz 组织；单位监护人。
                                \end{itemize}
                          \item 遗嘱指定。\zd{正在担任监护的父母}。尊重死者意思自治，效力具体情况具体分析。
                          \item 协议监护。防范风险，预先安排。开家庭会议，无顺序、人数限制。\zd{协议 \yjt 起诉 \yjt 指定监护}。
                          \item 指定监护。无顺序、人数限制。4 个单位，居委会、村委会、民政部门、法院。
                          \item 成年人意定监护。\zd{本质委托合同}。解决人口老龄化问题。
                          \item 委托监护。\zd{找人帮忙带孩子}。\al{留守儿童}。
                        \end{enumerate}
                  \item 职责。抚养、教育、保护。
                        \begin{enumerate}
                          \item 财产。不得随意处分，价值判断即可。
                          \item 人身。遗体捐献 = 父母 + 书面。
                        \end{enumerate}
                  \item 撤销与恢复。
                        \begin{enumerate}
                          \item 积极的侵权行为。\al{性侵害、出卖、虐待、遗弃、虐待、暴力伤害}。
                          \item 消极的侵权行为。\al{违法犯罪被关监狱，孩子无人看管}。
                          \item 向法院申请，\zd{民政局兜底}；3 费（抚养费、赡养费、抚育费）继续履行。
                          \item 恢复。\zd{故意犯罪的除外}。主体：只能是父母、子女。
                        \end{enumerate}
                  \item 终止。4 种。长大了，病好了；变傻了，或者一方死亡了。
                \end{enumerate}
          \item 宣告失踪。既不影响人身，也不影响财产关系，给失踪人的\zd{财产}找一个\zd{代管人}。参看民诉。
                \begin{enumerate}
                  \item 2 年。
                  \item 不申请宣告失踪，不影响其权利行使，义务履行。\al{甲乙夫妻借丙 100W 买房，后甲失踪}。
                \end{enumerate}
          \item \zd{宣告死亡}。
                \begin{enumerate}
                  \item 4 年；意外事件 2 年；无法生还随时。
                  \item \zd{利害关系人} \yjt 基层法院 \yjt 公告 1/1/3 \yjt \zd{宣告死亡} \yjt \zd{亡者归来}。
                  \item 法律后果。\zd{妻离子散，家破人亡}。
                        \begin{enumerate}
                          \item 妻离。婚姻关系自动解除。
                          \item 子散。死亡一方的父母有优先抚养的权利。
                          \item 家破。分遗产。
                          \item 人亡。视为自然死亡，丧失民事主体资格。
                        \end{enumerate}
                  \item 利害关系人顺序问题。立法的几度变更。
                        \begin{enumerate}
                          \item 无顺序。身份权由财产权来决定就很离谱。
                          \item 有顺序。4 个。配偶 > 父母 + 子女 > 其他近亲属 > 其他。bug \yjt 领空饷。
                          \item 无顺序。
                          \item \zd{无顺序但有条件}（当前）。
                                \begin{itemize}
                                  \item 第 1 组无条件 \yjt 父母、配偶、子女、丧偶儿媳/女婿。
                                  \item 第 2 组有条件 \yjt 其他近亲属 + 代位继承人。\sjt 均死亡 \hz 下落不明；不宣告\blh
                                  \item 第 3 组有条件 \yjt 债权人 + 债务人 + 合伙人。不宣告不足以保护其合法权益。
                                \end{itemize}
                        \end{enumerate}
                        % https://q.uiver.app/?q=WzAsNSxbMCwxLCLouqvku73mnYNcXFxc5LmZXFxcXOmFjeWBtuadgyJdLFsxLDEsIueUsiJdLFsyLDEsIuS5mVxcXFzvvIjlgLrmnYPkurogXFx5anQg6LSi5Lqn5p2D77yJIl0sWzEsMiwi5a2Q5aWzXFxcXOS6suWtkOadg++8iOebkeaKpO+8iVxcXFzouqvku73mnYMiXSxbMSwwLCLouqvku73mnYNcXFxc54i25q+NXFxcXOS6suWxnuadgyJdLFswLDEsIuWkq+Wmu+WFs+ezuyJdLFsxLDIsIjEwMFciXSxbMSwzXSxbMSw0XV0=
\[
  \begin{tikzcd}[mytest]
    & \txt{身份权\\父母\\亲属权} \\
    \txt{身份权\\乙\\配偶权} & \txt{甲} & \txt{乙\\（债权人 \yjt 财产权）} \\
    & \txt{子女\\亲子权（监护）\\身份权}
    \arrow["\txt{夫妻关系}", from=2-1, to=2-2]
    \arrow["\txt{100W}", from=2-2, to=2-3]
    \arrow[from=2-2, to=3-2]
    \arrow[from=2-2, to=1-2]
  \end{tikzcd}
\]

                  \item 亡者归来。妻离子散，家破人亡 \yjt 不一定不一定不一定不一定。
                  \item vs 宣告失踪。
                        \begin{enumerate}
                          \item 2 个独立的程序，没啥关系，没有顺序要求。
                          \item 一方失踪，一方死亡，如果满足宣告死亡条件 \yjt 宣告死亡。
                        \end{enumerate}
                \end{enumerate}
        \end{enumerate}
  \item 法人。依据法律创设的一种权利主体。\zd{公司}。合一主义，权利和行为能力同生同死，反面分离主义。
        \begin{enumerate}
          \item 参看商法。
          \item \zd{财产、人格、责任独立}。
          \item \zd{有限/无限责任，责任主体以其特定/全部财产承担责任}，内外有别。
          \item 始于成立，终于终止。设立中、存续中、清算中。
                \begin{enumerate}
                  \item 设立中。对外合同，一看结果，二看名义。
                        \begin{enumerate}
                          \item 未成立，\zd{视为合伙关系}（共同出资，共同经营，共享收益，共担风险）。
                          \item 成立，择一选择。
                        \end{enumerate}
                  \item 超越经营范围的法人。\zd{越权相对无效理论}，原则有效，\zd{禁}止、\zd{限}制、\zd{特}许经营无效。
                \end{enumerate}
          \item 分类。成立的目的和功能。
                \begin{enumerate}
                  \item 营利。非公司制，名称不带“有限责任”，“股份”；登记时成立。
                  \item 非营利。\zd{不分配利润}。公益 \hz 非公益。
                        \begin{enumerate}
                          \item 事业单位法人。\zd{公立的}学校、医院。
                          \item 社会团体法人。\al{律师协会；中国法学会}。
                          \item \zd{捐助法人}。\al{基金会（参照公司制度）；寺庙、道观；图书馆；博物馆}。
                        \end{enumerate}
                  \item 特别。
                \end{enumerate}
          \item 组成。权利、执行、监督机构。
          \item \zd{法定代表人}。\zd{代表权的限制，都是内部限制}。表见代表。
        \end{enumerate}
  \item 非法人组织。\zd{还是商事制度}。个人独资企业、合伙企业、不具备法人资格的专业服务机构（\al{分公司}）。\al{律所}。
        \begin{enumerate}
          \item 财产相对独立；补充的无限连带责任，因为空手套白狼。
          \item 分支机构。相对独立的民事主体资格，有程序而无实体的部分民事权利能力。\al{银行的分行、支行}。
        \end{enumerate}
\end{enumerate}
